\begin{itemize}
    \item Building Load contains : \textbf{gravity} and \textbf{lateral} Loads.
\end{itemize}
\subsection{Convert Area Loads to Line Load on Beams}
Slabs distribute their load to beam depending on slab span direction and support condition (Eurocode 2, Section on Slabs \cite{EN1992-1-1-2023(1)}; IStructE Manual, Chapter 4 \cite{IStructE2006(6)}).
\subsubsection*{The Tributary Width Method (IStructE Manual, Chapter 4 \cite{IStructE2006(6)}):}
For a \textbf{one-way slab}, load goes to the beams on the short sides:
\begin{center}
\begin{tcolorbox}[colback=white]
Beam Line Load (kN/m) = Floor Load (kN/m2) \(\times\) Tributary Width (m)
\end{tcolorbox}
\end{center}
\textbf{Tributary Width} = half the slab span on each side of the beam (IStructE Manual, Cl. 3.2 \cite{IStructE2006(6)}). \\[2mm]
\textbf{Example}: If Slab spans 5m between two beams, so each beam take \(5/2 = 2.5m\) tribuyary width. And if floor load is 10kN/m2, then beam line load is \( 10 \times 2.5 = 25 kN/m\) \\[2mm]
For a \textbf{two-way slab}, load distributes in both directions
(Eurocode 2, Section on Two-Way Slabs \cite{EN1992-1-1-2023(1)}; IStructE Manual, Chapter 5 \cite{IStructE2006(6)}):
\begin{itemize}
    \item Short span beam take triangular load
    \item Long span beam take trapezoidal load
\end{itemize}
\subsection{Apply ULS and SLS Load Combinations}
\begin{itemize}
    \item ULS (Ultimate Limit State) - for stength design
    \[\sum_{i} \gamma_{G,i} G_{k,i} + \gamma_{Q,1} Q_{k,1} + \sum_{j>1} \gamma_{Q,j} \varphi_{0.j} Q_{k,j}\]
    \item SLS (serviceability Limit State) - for deflection/crack check \\[2mm]
    Chearacteristic Combinations: \[\sum_{i}G_{k,i} + Q_{k,1} + \sum_{j>1} \varphi_{0,j} Q_{k,j} + (P_k)\]
    Quasi-permanent (for long-term deflection)
    \[\sum_{i}G_{k,i} + \sum_{j}\varphi_{2,j} Q_{k,j} + (P_k)\]
\end{itemize}
\subsection{Beam Support Conditions}
Support conditions determine the bending moment diagram shape and therefore where reinforcement must be placed
(Eurocode 2, Section on Structural Analysis \cite{EN1992-1-1-2023(1)}; IStructE Manual, Chapter 4 \cite{IStructE2006(6)}):
\begin{table}[H]
\centering
\caption{Beam support condition}
\resizebox{\linewidth}{!}{
\begin{tabular}{LLL}
\toprule
\textbf{Support Type}     & \textbf{Description}          & \textbf{Moment at Support} \\\midrule
\textbf{Simply supported} & Beam sits on wall/pad         & Zero moment at ends        \\
\textbf{Continuous beam}  & Multiple spans, monolithic    & Hogging moment at supports \\
\textbf{Fixed end}        & Fully restrained (rare in RC) & Full fixity moment         \\
\textbf{Partially fixed}  & Column/beam junction          & Partial moment transfer \\ \bottomrule
\end{tabular}}
\end{table}
\begin{table}[H]
\centering
\resizebox{\linewidth}{!}{
\begin{tabular}{LLLL}
\textbf{Support Type} & \textbf{Rotation Allowed} & \textbf{Moment at Support} & \textbf{Typical Structure} \\
Simply Supported      & Yes                       & 0                          & Beam on columns            \\
Continuous            & Small rotation            & Negative moment            & Multi-span floor beam      \\
Fixed                 & No rotation               & Maximum moment             & Cantilever                 \\
Partially Fixed       & Limited rotation          & Partial moment             & Beam-column joint
\end{tabular}}
\end{table}
\noindent In reinforced concrete buildings, beams are almost always continuous over columns because they are cast monolithically
(IStructE Manual, Chapter 4, Cl. 4.2 \cite{IStructE2006(6)}).
\subsection{Design Workflow}
\begin{itemize}
  \item \textbf{STEP 1:} Establish building layout and use

  \item \textbf{STEP 2:} Identify slab type (one-way / two-way)

  \item \textbf{STEP 3:} Gather \(G_k\), permanent loads \\
  (EN 1991-1-1:2002, Annex A, Table A.1 \cite{EN1991-1-1-2002(3)})

  \item \textbf{STEP 4:} Gather \(Q_k\), variable loads \\
  (EN 1991-1-1:2002, Cl. 6.3.1, Table 6.2 \cite{EN1991-1-1-2002(3)})

  \item \textbf{STEP 5:} Apply ULS combination

  \item \textbf{STEP 6:} Convert to beam line load

  \item \textbf{STEP 7:} Calculate bending moments and shear \\
  (IStructE Manual, Ch. 4 \cite{IStructE2006(6)}) \\
  (Eurocode 2, Structural Analysis \cite{EN1992-1-1-2023(1)})

  \item \textbf{STEP 8:} Size beam using \(L/d\) ratio \\
  (Eurocode 2, Deflection Section \cite{EN1992-1-1-2023(1)})

  \item \textbf{STEP 9:} Design reinforcement \\
  (Eurocode 2, ULS Design Section \cite{EN1992-1-1-2023(1)})

  \item \textbf{STEP 10:} Check SLS — deflection and cracking \\
  (EN 1990:2002, Cl. 6.5 \cite{EN1990-2002(5)}) \\
  (Eurocode 2, SLS Section \cite{EN1992-1-1-2023(1)})

  \item \textbf{STEP 11:} Detail and produce drawings \\
  (Eurocode 2, Detailing Section \cite{EN1992-1-1-2023(1)}) \\
  (IStructE Manual, Ch. 4 \cite{IStructE2006(6)})
\end{itemize}
\subsection{The Key Distinction: what are we Designing}
According to
(IStructE Manual, Chapter 3 \cite{IStructE2006(6)}):
\begin{table}[H]
\centering
\begin{tabular}{LL}
\toprule
\textbf{Element} & \textbf{How Load Works}               \\\midrule
Beam             & Carries only load from its own floor  \\
Column           & Carries load from all floors above it \\
Foundation       & Carries load from the entire building \\ \bottomrule
\end{tabular}
\caption{How structure work}
\end{table}
\subsection{The Unit Journey}
\begin{table}[H]
\resizebox{\linewidth}{!}{
\begin{tabular}{LLLL}
\toprule
\textbf{Unit} & \textbf{Name} & \textbf{Used For}        & \textbf{How We Get It}    \\ \midrule
kN/m³         & Density       & Material property        & Given in \cite{EN1991-1-1-2002(3)} Table A.1 \\
kN/m²         & Area load     & Slab design              & Density \(\times\) thickness        \\
kN/m          & Line load     & Beam design              & kN/m² \(\times\) tributary width    \\
kN            & Point force   & Column/foundation design & kN/m \(\times\) span \(\div\) 2            \\ \bottomrule
\end{tabular}}
\caption{Unit of load to consider in design}
\end{table}